\chapter*{Lời mở đầu}\addcontentsline{toc}{chapter}{Lời mở đầu}

Xổ số kiến thiết là một ví dụ điển hình của quá trình ngẫu nhiên trong thực tế, 
trong đó kết quả của mỗi kỳ quay về nguyên tắc được kỳ vọng là độc lập và không 
thể dự đoán một cách có hệ thống từ các kết quả đã quan sát trong quá khứ. Tuy 
nhiên, khi một lượng lớn dữ liệu lịch sử được tích lũy qua thời gian, một câu hỏi 
tự nhiên được đặt ra là liệu các kết quả xổ số thực tế có thực sự thể hiện đầy đủ 
các đặc trưng của một quá trình ngẫu nhiên hay vẫn tồn tại những sai lệch thống kê, 
sự phụ thuộc theo thời gian hoặc các cấu trúc có thể khai thác được.
Trên thế giới đã có những nghiên cứu và công bố về học thuyết xổ số, tuy nhiên không có 

Xuất phát từ vấn đề trên, bài nghiên cứu này phân tích dữ liệu kết quả xổ số kiến thiết dựa trên các kiến thức của môn học Các phương pháp ngẫu nhiên và ứng dụng, đồng thời cũng áp dụng các mô hình máy học hiện đại. Mục tiêu của nghiên cứu là kiểm chứng giả thuyết bằng dữ liệu thực nghiệm. Nghiên cứu 
hướng tới làm rõ ba câu hỏi chính: \textit{(i)} kết quả xổ số có phù hợp với các đặc trưng kỳ 
vọng của một quá trình ngẫu nhiên hay không; \textit{(ii)} các mô hình xác suất và học máy có 
khai thác được tín hiệu dự báo ổn định ngoài mẫu hay không; \textit{(iii)} và nếu tồn tại tín hiệu, 
mức cải thiện đó có đủ để tạo ra một chiến lược có hiệu quả kinh tế sau khi tính 
đầy đủ chi phí tham gia và cơ cấu giải thưởng hay không.